% Stage 1 presentation: MediSlot pentest project overview (max 5 slides).
% Build: pdflatex stage-1-presentation.tex (run twice).
\documentclass[aspectratio=169,11pt]{beamer}

\usepackage[utf8]{inputenc}
\usepackage[T1]{fontenc}
\usepackage{lmodern}
\usepackage[polish]{babel}
\usepackage{booktabs}
\usepackage{tikz}
\usetikzlibrary{arrows.meta,positioning}

\usetheme{Madrid}
\usecolortheme{seahorse}
\setbeamertemplate{navigation symbols}{}
\setbeamertemplate{caption}[numbered]

\definecolor{crit}{HTML}{B00020}
\definecolor{high}{HTML}{C8651B}

\title[MediSlot — pentest]{Test penetracyjny środowiska MediSlot}
\subtitle{NovaMed Sp. z o.o. — etap I (SoW + Threat Model)}
\author[SecureTrio]{Zespół SecureTrio: Kamil \and Bartek \and Kuba}
\institute[AGH]{Testy penetracyjne — projekt zaliczeniowy, AGH}
\date{24.10.2026}

\begin{document}

% ---------- Slajd 1: tytuł ----------
\begin{frame}[plain]
  \titlepage
  \begin{center}
    \footnotesize Lab lokalny (Docker Compose) \textbullet{} dane syntetyczne \textbullet{} model grey-box
  \end{center}
\end{frame}

% ---------- Slajd 2: co testujemy + zakres ----------
\begin{frame}{Co testujemy: MediSlot i stack labowy}
  \textbf{MediSlot} — portal przychodni: rejestracja wizyt, dokumentacja pacjentów (syntetyczna), 4 role (pacjent, lekarz, recepcja, admin).

  \begin{columns}[T]
    \begin{column}{0.52\textwidth}
      \begin{block}{Technologie (7 komponentów, Docker Compose)}
        \footnotesize
        \begin{itemize}
          \item \textbf{Nginx} — reverse proxy
          \item \textbf{PHP-FPM} — aplikacja MediSlot
          \item \textbf{MySQL} — baza danych
          \item \textbf{Redis} — magazyn sesji
          \item \textbf{MinIO} — załączniki (S3)
          \item \textbf{Mailhog} — poczta labowa
          \item \textbf{Debian} — warstwa OS kontenerów
        \end{itemize}
      \end{block}
    \end{column}
    \begin{column}{0.48\textwidth}
      \begin{block}{Zakres}
        \footnotesize
        \textbf{In scope:} aplikacja, API HTTP, konfiguracja usług, ruch w sieci Compose (\texttt{localhost}), kontrola dostępu między rolami, łańcuchy do RCE.\\[0.4em]
        \textbf{Out of scope:} host Windows / Docker Desktop, sieci poza labem, prawdziwe dane/PII, DoS hosta, infrastruktura AGH.
      \end{block}
    \end{column}
  \end{columns}
\end{frame}

% ---------- Slajd 3: plan, metodyka, narzędzia ----------
\begin{frame}{Plan, metodyka i narzędzia}
  \begin{columns}[T]
    \begin{column}{0.54\textwidth}
      \begin{block}{Plan (etapy projektu)}
        \footnotesize
        \begin{enumerate}
          \item \textbf{Etap I} — SoW + Threat Model (+ ta prezentacja)
          \item \textbf{Etap II} — testy: discovery $\rightarrow$ atak $\rightarrow$ dowody
          \item \textbf{Etap III} — raport końcowy + re-testy
        \end{enumerate}
      \end{block}
      \begin{block}{Metodyka}
        \footnotesize
        PTES / OWASP Testing Guide + mapowanie na \textbf{STRIDE} z threat modelu. Każdy finding: opis, reprodukcja, dowód, CWE, CVSS, CVE, remediacja.
      \end{block}
    \end{column}
    \begin{column}{0.46\textwidth}
      \begin{block}{Narzędzia}
        \footnotesize
        \begin{itemize}
          \item Proxy: \textbf{Burp} / \textbf{ZAP}
          \item \textbf{sqlmap} — SQLi
          \item \textbf{Trivy} — skan obrazów / CVE
          \item \textbf{Lynis} — audyt OS
          \item \textbf{Nuclei} + własne skrypty HTTP
        \end{itemize}
      \end{block}
      \begin{block}{Docelowo}
        \footnotesize $\sim$15–20 findingów: kilka Critical/High, większość Medium, reszta Low/Info.
      \end{block}
    \end{column}
  \end{columns}
\end{frame}

% ---------- Slajd 4: podatność 2026 + top ryzyka ----------
\begin{frame}{Podatność z 2026 i top ryzyka}
  \begin{block}{Główny tor Critical — podatność 2026}
    \footnotesize
    Celowo podatna wersja \textbf{PHP-FPM} z \textbf{CVE z 2026 r.} prowadząca do \textbf{RCE} (pełne przejęcie kontenera aplikacji). Discovery przez fingerprinting / skan obrazu (Trivy), potwierdzenie eksploatacją w granicach RoE.
  \end{block}

  \centering
  \footnotesize
  \begin{tabular}{@{}llccl@{}}
    \toprule
    \textbf{ID} & \textbf{Zagrożenie} & \textbf{P} & \textbf{W} & \textbf{Priorytet} \\
    \midrule
    R1 & RCE przez podatny PHP (CVE 2026) & 4 & 5 & \textcolor{crit}{Critical — RoE} \\
    R2 & RCE/łańcuch przez upload (MinIO) & 3 & 5 & \textcolor{crit}{Critical} \\
    R3 & IDOR na karty / notatki / wizyty & 5 & 4 & \textcolor{high}{High} \\
    R4 & Vertical privilege escalation (rola) & 4 & 5 & \textcolor{crit}{Critical/High} \\
    R5 & Stored XSS $\rightarrow$ przejęcie sesji & 4 & 4 & High/Medium \\
    R6 & Słaby reset hasła + Mailhog & 4 & 4 & High/Medium \\
    \bottomrule
  \end{tabular}

  \vspace{0.4em}
  {\footnotesize P = prawdopodobieństwo, W = wpływ (skala 1–5).}
\end{frame}

% ---------- Slajd 5: scenariusz testera vs obrona (SOC) + RoE ----------
\begin{frame}{Scenariusz testera vs obrona (SOC)}
  \begin{columns}[T]
    \begin{column}{0.5\textwidth}
      \begin{block}{Scenariusz testera (atak)}
        \footnotesize
        \begin{enumerate}
          \item Fingerprinting PHP / skan obrazów
          \item Kontrola dostępu / IDOR na wizytach i kartach
          \item Upload / Stored XSS / reset hasła
          \item Eksploatacja \textbf{CVE PHP} jako tor Critical (RCE)
        \end{enumerate}
      \end{block}
    \end{column}
    \begin{column}{0.5\textwidth}
      \begin{block}{Scenariusz obrony (wiedza SOC)}
        \footnotesize
        \begin{itemize}
          \item Logi aplikacji: logowania, błędy, upload
          \item Access/error logi Nginx
          \item Skan obrazów (Trivy) — alert przy krytycznym CVE
          \item Możliwość wyłączenia konta; brak WAF w labie
        \end{itemize}
      \end{block}
    \end{column}
  \end{columns}

  \begin{block}{Procedura RoE przy ustaleniu Critical}
    \footnotesize
    Natychmiastowy kontakt z RoE (Bartek) $\rightarrow$ wstrzymanie eksploatacji wektora $\rightarrow$ notatka (timing, wektor, wpływ) $\rightarrow$ dyspozycja (kontynuuj / wstrzymaj / zbierz logi). Re-test findingów Critical i High po patchu.
  \end{block}

  \centering\footnotesize Dziękujemy — pytania? \quad\textbullet\quad Dokumenty: SoW + Threat Model (etap I).
\end{frame}

\end{document}
